% 23_rebuttals_and_ac_comments.tex — Writing the Author Response: Rebuttals and the AC Comment

\section{Writing the Author Response: Rebuttals and the AC Comment}
\label{sec:rebuttals_and_ac_comments}

The publication lifecycle covered so far ends at submission.
The step immediately after it---the author response, or rebuttal---receives almost no explicit training, and the feedback loop is slow: an author writes one or two rebuttals per year and never observes the counterfactual.
This section records the heuristics we applied during a recent NeurIPS position-paper reviewing cycle, refined by advisor and co-author review of our own draft responses.
We state them as heuristics from a single cycle rather than as established findings; venues differ, and area chairs differ within a venue.
The organising claim is that the response has one primary reader, and it is not the reviewer who wrote the harshest sentence.

\subsection*{Priority: The Response Is Written for the Area Chair}

Reviewers rarely revise scores.
An author can answer a weakness cleanly, receive an acknowledgement that the concern is addressed, and see the numeric score remain unchanged; the reviewer has already committed a position to the record, and updating it carries a cost.
The area chair (AC) is differently situated.
The AC must produce a decision from the entire record---every review, every response, the reviewer discussion, and their own reading---and is the only participant who sees all threads side by side.

Two consequences follow.
First, \textbf{answer the AC metareview before addressing individual reviews}, while attention is still fresh.
If the metareview names three concerns, those three concerns are the paper's operative weaknesses for decision purposes, irrespective of what any single reviewer chose to emphasise.
Second, \textbf{cross-thread consistency becomes load-bearing rather than cosmetic}.
Describing an effect as small in one thread and as clear in another is visible only to the reader who holds both, which is precisely the reader who decides.

\textbf{After the metareview, bias effort by score: lowest first, then the middling reviews, and leave a strong accept largely alone.}
The rationale is again about the AC's position.
A uniform set of accepts is straightforward to defend in committee; one strong accept adjacent to a reject is not, and a single low score obliges the AC to justify the outcome.
The corollary is less obvious and worth stating: a strong accept can be lost during the rebuttal, not by pushing back, but by conceding so extensively across the response that the reviewer re-examines their own generosity.
Concession density deserves attention on the reviews that are already favourable.

\subsection*{Whether to Concede}

\textbf{Never concede blindly.}
Concede a genuine error, or a limitation that the evidence in hand actually establishes.
Otherwise clarify, or point to material the paper already contains, which is the correct answer more often than authors under time pressure assume.

\textbf{Verify a technical premise before accepting it.}
Reviewer premises are sometimes factually wrong.
A representative case: a review asserts that the comparison used a weaker baseline, when the two methods are the same algorithm published under different names.
Conceding that premise commits a false statement to the record in a place the AC will read.

One pattern is easy to misclassify as an attack.
When a reviewer proposes an unmeasured property as an alternative explanation for a reported effect, and the paper's position is that such properties should be measured and reported, the reviewer has supplied support for the position rather than a refutation of it.
The appropriate response credits the reviewer, agrees, and states what the alternative explanation implies for practice.

Where the honest answer lies between the reviewer's reading and the paper's, \textbf{prefer calibrated refinement to dramatic retreat}.
``A clear change in observed behaviour, small in magnitude'' is defensible under scrutiny.
``We withdraw the claim'' invites the AC to ask which remaining claims are similarly soft.

\subsection*{How to Concede}

The wording of a concession carries as much weight as the decision to make one, and this is where drafts most often fail internal review.

Formulations such as ``we concede two points outright,'' ``we withdraw,'' and ``this is our error'' surrender more than the reviewer requested and make the paper sound broken to a reader who is skimming.
A co-author will object to such a sentence before reaching the substance beneath it, which is a reliable predictor of how an AC will read it.

\textbf{Concede the specific thing, never the category.}
``That word claims more than the numbers support, and we will change it'' is a bounded wording fix.
``The claim is unsupported'' is an admission about the paper's evidence, and the reviewer did not ask for it.

\textbf{Never volunteer self-criticism harsher than the reviewer's own framing.}
A reviewer writing that two constructs are ``connected but not the same'' has made a request for precision, not an accusation of error, and should be answered as such.
Intensifiers on concessions cost credibility and buy nothing: ``outright,'' ``without qualification,'' ``entirely,'' ``completely,'' ``we were wrong.''
The unmodified verb is stronger.
Nor should an author direct a reviewer to the paper's weakest point.
Writing ``this is the sharpest gap in the paper'' happens, usually out of a sense of candour; the reviewer has already scored the submission and does not need the assistance.

Two rules do most of the work in practice.
\begin{enumerate}
    \item \textbf{Every concession names what changes and where.}
    ``We will revise the phrasing in \S4.2 and add the caveat to the caption of Table 3.''
    A concession with no stated consequence reads as evasion even when it is sincere.
    \item \textbf{Every concession appears alongside what survives it}, in the same paragraph.
    Give up the overclaiming word, then state the finding the word was overclaiming, which remains intact.
\end{enumerate}

Close the response on what stands, not on a tally of what changed.

\subsection*{Evidence}

\textbf{Every number in a response must be traceable to a table, section, or line the reviewer can open.}
A practical test: if a co-author cannot locate the number, a reviewer certainly cannot.

\textbf{Prefer facts already present in the submitted PDF to anything generated inside the rebuttal window.}
The strongest available answers are frequently already in the paper's own appendix tables, unread because appendices are long and reviewers are busy.
Directing a reviewer to an existing table is more persuasive than a hurried new experiment, and considerably cheaper.

When new analysis is unavoidable, label it accurately.
\textbf{A post hoc comparison is a sensitivity check, not a controlled experiment}, and describing it as the latter is the kind of imprecision an attentive reviewer will notice and remember.
State $n$.
When the evidence is that three conditions fell in the predicted order, state that a perfect ordering of three arises by chance with probability $1/6$; it is better said by the authors than by the reviewer.
Keep association and causation distinct, and where only association is available, name the design that would identify the causal claim.
\textbf{Pre-empt what a careful reviewer would find}, including the outlier row in the very table being cited.

Two further constraints:
\begin{itemize}
    \item \textbf{Results, not promises.}
    Reviewers discount promised experiments heavily, and the discount is warranted.
    \item \textbf{Do not cite an adjacent unpublished or parallel paper of your own.}
    It compromises double-blind review, and it invites a critique that cannot be answered inside the response: why is this two papers if the load-bearing evidence lives in the other one?
    Absorb the scope gap instead.
\end{itemize}

\subsection*{Construction}

\textbf{Answer the question actually asked, not the adjacent question that is easier to answer.}
``What is the consequence for an experimentalist?'' requires concrete actions; restating the mechanism is a non-answer, and a reviewer who asked a practical question will register it as one.

\textbf{Follow each reviewer's own numbering}---W1, W2, Q1---so that every item is visibly discharged.
This serves the reviewer and, more importantly, the AC scanning for unanswered items.
Thank reviewers for strengths specific to their review, briefly; generic praise of a thin review is legible as such to the AC reading both documents.
When a reviewer supplies a useful framing and asks that it be made explicit, adopt it, credit them in a clause, and add something beyond it.
Pure echo is a thank-you note, and it forfeits an opportunity to show that the framing extends further than the reviewer suggested.

Smaller mechanical points that erode under time pressure: use third person for all self-citation, without exception; do not pad to the character limit, since length is not the objective; avoid absolutes such as ``solves selection bias,'' ``isolates $X$,'' ``$X$ causes $Y$,'' and ``method $M$ wins''; keep numbers and claims consistent across all threads.
Ask what would raise the reviewer's score.
Never ask for a higher score.

\subsection*{Voice: The Response Must Not Read as Machine-Written}

Reviewers penalise author responses that read as language-model output, and some tracks prohibit AI-written responses outright.
Whatever one's view of the policy, it is the operative incentive.
The target register is a tired expert writing carefully.

The reliable tells, roughly in descending order of how loudly they announce themselves:
\begin{itemize}
    \item \textbf{Uniform sentence length.}
    A run of same-length balanced sentences is the strongest single signal.
    \item \textbf{Recurring frames}: ``not $X$, but $Y$,'' ``it is worth noting that,'' ``this is precisely why,'' ``we would like to emphasise,'' ``delve,'' ``leverage,'' and ``robust'' used as filler.
    \item \textbf{Em-dash saturation.}
    Roughly one per section; commas, colons, and full stops remain available.
    \item \textbf{Triads used for rhythm rather than content}, as in ``clear, actionable, and testable.''
    \item \textbf{Consecutive paragraphs with identical openings}, such as two paragraphs beginning ``We agree.''
    \item \textbf{Symmetrical section openings.}
    Human writing is uneven, and one section being blunter than the others reads as authored.
    \item \textbf{Abstractions where a number would serve.}
    ``Table 7 gives $-0.105$'' is stronger than ``the empirical evidence indicates.''
\end{itemize}

Contractions in moderation help; a perfectly formal register maintained throughout reads as generated.
Read the response aloud before posting, which catches more than any checklist, including the one in Appendix~\ref{appendix:rebuttal_checklist}.

\subsection*{The AC Comment, Written Last}

The comment to the area chair is drafted last, because it summarises the reviewer threads and cannot be written before those threads are settled.
\textbf{Mirror the AC's stated priorities, in the AC's order, with the AC's numbering.}
Three concerns in the metareview become three parts in that sequence.
Then close with an enumerated list of manuscript changes and the location of each.
That list is the most useful paragraph in the entire response, because it is the item the AC can point to when writing the decision.

\subsection*{Limits of This Advice}

This is one cycle, one track, one paper, and the specifics should be treated accordingly.
The claim we would defend most firmly is that the AC is the primary audience, which follows from the structure of the process rather than from our experience of it.
The claim we would defend least is the concession-density argument---that a strong accept can be lost by over-conceding---which we believe on the basis of advisor and co-author judgement and cannot establish from a single cycle.
Venues also vary in how closely area chairs read author responses, and an author cannot determine which case they are in until the decision is released.
What we expect to generalise is the underlying reorientation: write for the participant who must defend a decision, rather than for the reviewer who caused the most irritation.
